\section{Coherence, Selves, and Mind}
\label{sec:mind}

A self is not a special vibe. It is a coherent, self-maintaining pattern of vibes and notes, and everything we call mind is what such a pattern does.

\subsection{Nesting and the combination problem}

When a set of vibes is bound by strong mutual notes and shares a tone, it behaves as a single higher vibe. It has one collective tone, projects one collective will, and can be noted as a unit. This is coarse-graining, or renormalization, the established mathematics of how many small units act as one larger unit. It is the model's answer to the combination problem, the question of how micro-experiences make a mind. A parent mind is a coherent block of child vibes, the same stuff at a larger scale. The tower of scales, from vibe to particle to cell to organism to world, is renormalization iterated.

\subsection{Coherence and the threshold of agency}

Coherence measures how tightly a block holds together, the strength of its internal notes combined with the agreement of its tones. As internal coupling rises past a critical value, the block crosses a phase transition. Below it the block merely reacts. Above it a single collective mode governs it, and it acts as a unified agent. So agency above a threshold is the onset of an ordered phase, a computable transition rather than a vague upgrade.

\subsection{The self is a Markov blanket}

A self is the internal vibes of a coherent block, separated from the external vibes by a blanket of notes, the notes through which the inside senses (fill) and acts (will) on the outside. Statistically this is a Markov blanket, a boundary that renders the internal states conditionally independent of the external. Self and other are therefore not metaphysical absolutes but a structural boundary in the note-field, dynamic and redrawn each beat. What is other can become self by being bound in, and self can become other by dissolving.

A natural minimal picture is a candle flame or a whirlpool in a stream. It is a pattern, not a substance. The matter passing through it is never the same twice, yet it holds a stable shape with a definite boundary, maintained against the flow by drawing in fuel and exporting heat. That is a self in miniature, a coherent pattern with a blanket boundary. The self is the form, not the stuff.

\subsection{Mind as the minimization of surprise}

A self that persists must keep its blanket intact against a changing environment. The principle that captures this is active inference (Friston, 2010). The self acts and perceives so as to minimize its free energy, its surprise, the mismatch between what it predicts and what it notes. Perception is updating internal tones to fit what is noted. Action is projecting will to change what is noted so it fits the model. Learning is the slow reshaping of note strengths toward a better model. So ``minimize suffering'' (Section~\ref{sec:energy}) and ``minimize free energy'' are the same imperative, and it generalizes mere tone-alignment. A self does not only smooth its tones, it predicts and acts to confirm its predictions, which is why it stays differentiated rather than dissolving into the environment.

\subsection{Will, choice, and learning}

Base will is a vibe's tone turned outward, mechanical and without purpose. Volition is not primitive. It emerges as base will aggregated over a coherent self and steered by the free-energy gradient. The self projects more strongly in the direction its model predicts will lower its suffering. Purpose is base will plus coherence plus the gradient. Because attention is structural (Section~\ref{sec:dynamics}), the self cannot steer by instantly re-attending. It steers by where it projects its will and by slowly reshaping its note-structure.

A coherent self does not flip instantly when its inputs change. Its internal vibes must re-reach internal agreement, which takes a number of beats, the settling time of its internal dynamics. That interval, experienced from inside, is deliberation. From outside it is a deterministic relaxation. Both are true. Why it still feels open is computational irreducibility. There is no shortcut to the outcome faster than running the self, so no observer can know the choice before it settles. Learning and karma are the same slow process. Notes between vibes that keep agreeing strengthen, carving the attractor landscape that is memory, and actions leave persistent note-changes that propagate until they balance.
